\section{Functional change and the contradiction}\label{sec:contradiction}

We apply the preceding results to the buffered intervals of
Proposition~\ref{prop:buffered-intervals}.  The interval is
$I=[-\theta,1+\theta]$ in its rescaled time coordinate.  The central
interval $[0,1]$ has original duration $d_i\to0$, whereas the full
buffered interval has original duration $(1+2\theta)d_i$.
Put $\epsilon_i=\sqrt{d_i}$.  We suppress the index $i$ when no limit
is being taken.  Thus $|R|\le C\epsilon_i^2$,
the reference root length $a=a_i$ tends to zero, and the root scale
defined in Proposition~\ref{prop:root-scale} is comparable to $a$
throughout $I$.  Its values at $0$ and $1$ differ by a fixed factor.
The centers $x_t$, cutoff $\eta$, and energy $K$ are those of
Equation~\eqref{eq:root-energy}; in particular
\begin{equation}\label{final:energy-recall}
 a^4 K^2=\int_I\eta(t)^2
       \int_{B_t(x_t,N_1a)}|\Ric|^2\,d\mu_{g(t)}\,dt,
 \qquad \min_{[0,1]}\eta>0.
\end{equation}
The maximal function $H_*$ may use a larger fixed ball than the one in
Equation~\eqref{final:energy-recall}, as permitted in
Equation~\eqref{eq:maximal-ricci}.

First fix the small collar tolerance.  The estimates below hold for
every sufficiently large fixed dyadic number $N\gg N_1$, with
$b=Na$.  All limits $i\to\infty$ are taken with $N$ fixed.
The constants multiplying $N^{-2\alpha}K^2$ will be independent of
large $N$; constants multiplying a quantity denoted $o_i(1)$ may depend
on $N$.  At the final absorption step we choose one such $N$ large
enough to absorb that explicit loss, then take the sequence limit.
No subsequent limit in $N$ is needed.

\subsection{Material collars and time patches}

\begin{lemma}[Position of a controlled material collar]
\label{final:collar-position}
Choose sufficiently small $\sigma_0>0$ and sufficiently large fixed
$N$.  At a time near $\supp\eta$, choose an exterior annular chart at
scale $b$ with domain, in collar units,
$\{1/8<|x|<16\}/\Gamma$, retaining slightly larger boundary margins
when necessary.  Keep its map into the material manifold and its group
fixed.  For all sufficiently late $i$, throughout any interval on
which its metric stays within $\sigma_0$ of the flat metric through the
required fixed derivative order, the following assertions hold.
Every closed shortened collar lies at distance comparable to $b$
from the current center $x_t$; its middle bands cover the corresponding
shortened distance annuli.  The component on the designated inner side
of a middle sphere contains $B_t(x_t,N_1a)$ and is contained in
$B_t(x_t,Cb)$, with fixed constants.
\end{lemma}

\begin{proof}
The initial chart has these properties by the exterior annular
convergence in Proposition~\ref{prop:exterior-ricci}.  In particular,
a middle sphere separates: every path from a sufficiently small ball
to outside a sufficiently large comparable ball crosses the chart's
radial band and hence the sphere.  Its two sides belong to different
components.  A connected two-sided embedded closed hypersurface in a
connected manifold has at most two complementary components: every
component reaches a tubular neighborhood of the hypersurface, and
each of its two connected sides meets only one component.  Thus the
designated inner component is unambiguous and is fixed as a material
subset while the chart is kept fixed.

We justify the uniform positional assertion; closeness of the metric
in a chart alone would not imply it.  If the assertion failed for
arbitrarily small tolerances and arbitrarily large $N$, choose a
countersequence with $\sigma_0\to0$, $N\to\infty$, and indices as late
as necessary for each choice.  The chart lies in the material
neighborhood of the same isolated point of the limiting space $Z$.
At scale $b$, based at an interior chart point, Theorem~\ref{thm:degeneration}
gives a complete flat quotient limit: a nonflat limit at this scale,
which is much larger than $a$, would contradict the outermost-scale
property in Proposition~\ref{prop:root-scale}.  On each compact
subannulus the controlled chart supplies curvature bounds on balls
with a definite margin.  Indeed a path leaving the chart must first
traverse that margin in its controlled metric.  Smooth compactness
therefore extracts the chart maps locally into the regular part of
the flat limit; metric and connection transformation formulas give
convergence of the maps and a limiting local flat isometry.

The groups can be fixed after passage to a subsequence.  Their orders
are bounded by noncollapse, there are finitely many abstract groups
of bounded order, and their orthogonal representations have convergent
subsequences.  Close orthogonal representations of a fixed finite
group are conjugate: averaging an approximate intertwiner over the
group makes it an exact invertible intertwiner, and its polar factor
is orthogonal and still intertwines.  The source group is nontrivial
and acts freely on the sphere.  The target flat limit is a global
Euclidean quotient with at most one exceptional point, by
Theorem~\ref{thm:degeneration}.  Lift the local isometry to Euclidean
covers of regular annuli.  It is the restriction of a rigid motion
$x\mapsto Ax+c$, since its differential is parallel.  Equivariance
conjugates the source group into the target group.  Its translation
part fixes $c$ under that nontrivial free spherical action, so $c=0$
after centering the target quotient.  In particular the target group
is nontrivial and the image is a complete centered open annulus.

There is no multiplicity in that image.  Otherwise two distinct
interior source points would approach one target point.  In the
approximations they could then be joined by a path of length $o(b)$.
A path contained in the chart cannot do this by its controlled
metric, and a path leaving it costs a fixed positive multiple of $b$
before reaching the boundary.  Both alternatives are impossible.
The target tip is represented by points $z_i$ with $l(z_i)/b_i\to0$.
It must lie at distance $o(b_i)$ from the current cluster center.
Otherwise, on a subsequence, put $h_i=d_t(z_i,x_t)\ge cb_i$; the
case $h_i/b_i\to\infty$ is allowed.  The chart's material points
approach the same point of $Z$ by their initial positioning and
uniform convergence of the original distances.  The same holds
for $z_i$, which is within $O(b_i)$ of an interior chart point,
and for $x_t$.  Thus the original length corresponding to $h_i$
tends to zero.  At this separation scale,
\[
 \frac{l(z_i)}{h_i}\longrightarrow0,\qquad
 \frac{l(x_t)}{h_i}\le C\frac{a_i}{h_i}
                    \le\frac{C}{cN_i}\longrightarrow0.
\]
Here the maximizing-center detection in
Proposition~\ref{prop:root-scale} gives $l(x_t)\asymp a_i$;
this placement countersequence has $N_i\to\infty$.
The separation-scale limit has two genuine tips at distance one.
It is therefore nonflat, since a flat limit has at most one tip.
Its scale satisfies $h_i/a_i\ge cN_i\to\infty$, contradicting
the outermost-scale property in Proposition~\ref{prop:root-scale}.

These facts imply coverage in the approximations.  A point whose
limit belongs to the interior image has distance $o(b)$ from an image
point with a fixed chart margin; a path realizing that distance
cannot leave the image, so the point itself lies in it.  Boundary
positions follow by approaching the boundary from interior points
and using the controlled lengths.  Consequently all shortened
radial bands have their claimed position and coverage, to any
prescribed fixed accuracy.  Crossing such bands shows that the
geometric inward component contains the small ball and is contained
in the larger $Cb$ ball.  To identify it with the designated material
component, fix a regular material point $p$ at positive $d_T$-distance
from this isolated cluster.  At the chart's initial time $p$ belongs
to the outward component.  Uniform convergence of original distances
keeps $p$ outside the current $Cb$ ball at every time under consideration.
Thus it still belongs to the geometric outward component.  The
embedded sphere and its two complementary material components were
kept fixed, so the other component remains the designated inward one.
This contradicts the countersequence and proves the lemma with fixed
choices of tolerance and $N$, followed by sufficiently late indices.
\end{proof}

At any individual time in such a patch, complete exterior coordinates
can be chosen to agree with the fixed chart on $C_1=\{1\le|x|\le2\}/\Gamma$.
Here and below an outward fixed rescaling of the chosen middle collar
is harmless.  To see the agreement, take the single neck coordinates
from Equation~\eqref{eq:neck-decay} out to a sufficiently small fixed
radius $r_*<r_0$ in interval units.  On the roomy overlap their
transition with the fixed chart has metric distortion
$O(\sigma_0+N^{-\mu})$ and controlled connection difference.  Lifting
the overlap, integration of the connection transformation makes
this transition close to a rigid equivariant map, through all needed
orders.  Its linear part can be adjusted to be orthogonal by polar
factorization.  The group and whole-band identifications are those
proved in Lemma~\ref{final:collar-position}.  After that rigid
alignment the maps are close to identity, so a radial cutoff
interpolation on a shorter overlap has invertible derivative and
patches the coordinate systems.  This is precisely the overlap
construction of the neck charts, now performed only in a regular
fixed-width band.  Decreasing $\mu>0$ if needed, the shell weights in
Theorem~\ref{thm:collar-extension} can consequently be chosen as
\begin{equation}\label{final:collar-weights}
 \theta_r\le C\bigl[(\sigma_0+N^{-\mu})r^{-\mu}
                    +\delta_i(rb)^\mu\bigr],
 \qquad 1\le r\le4R,\qquad Rb\asymp r_*.
\end{equation}
They include the collar norm times $r^{-q}$ after decreasing $\mu<q$.
The dyadic sum of their first term is small with the fixed choices;
the second sum tends to zero because $Rb\asymp r_*$.
These coordinates outside $C_1$ are used only to estimate the error.
They are not differentiated in time.

Let $B_1(t)$ be an upper bound for $|\Ric|_{k+4,b}$ on a fixed ample
band of true distance radii $[c'b,C'b]$ around $x_t$, in interval
units.  Suprema over all admissible center choices can be used to
make this bound independent of a choice of chart.  It is measurable
and locally bounded for every $i$.  Define an increasing absolutely
continuous time coordinate by
\begin{equation}\label{final:time-coordinate}
 u(t)=\int_{t_*}^{t}(1+B_1(s))\,ds .
\end{equation}
Start each chart with a tolerance strictly better than $\sigma_0$,
as is possible by first taking $N$ large.  As long as the chart
remains within $\sigma_0$, Lemma~\ref{final:collar-position} places
its entire closed domain in the band defining $B_1$.  The flow equation
$\partial_tg=-2\Ric$, together with conversion between covariant and
coordinate derivatives in that controlled chart, then gives
\[
 \|g(t)-g(s)\|_{C^{k+3}(\text{fixed collar})}
       \le C\int_s^t B_1(v)\,dv .
\]
A first-exit argument proves validity for a fixed sufficiently small
$u$-width.  This holds in both time directions, using the absolute
value of the integral.

On an interval with room around $\supp\eta$, choose regularly spaced
centers in the $u$ coordinate, subordinate smooth nonnegative bumps
of that fixed width, and normalize their sum.  Pulling them back gives
an absolutely continuous partition $\{\chi_j\}$ with
\begin{equation}\label{final:partition}
 \sum_j\chi_j=1\quad\hbox{near }\supp\eta,\qquad
 \sum_j|\chi_j'(t)|\le C(1+B_1(t))\quad\hbox{a.e.}
\end{equation}
Its overlap is bounded by a fixed number.  For each $i$ it is finite
on the compact interval in use; no bound on the number of members
uniform in $i$ is asserted or needed.

On patch $j$, retain the true metric on the material inner component
and apply Theorem~\ref{thm:collar-extension} in its fixed collar.
Interpolate as in that theorem.  Denote the resulting metric on the
fixed manifold with one end by $\hat g_j(t)$, and its functional in
interval length units by $E_j(t)$, using
Equation~\eqref{eq:energy-functional}.  Different patches are allowed
to have different extensions and different functional values.

\subsection{The lower bound for functional variation}

\begin{lemma}[The time-integrated collar error]
\label{final:collar-bound}
The residual in Equation~\eqref{eq:collar-residual} for every active
patch is bounded by a common nonnegative function $m(t)$ such that
\begin{equation}\label{eq:collar-time-bound}
 \sup_{\supp\eta}m=o_i(1),\qquad
 b^{-2}\|\eta m\|_{L^2(I)}+\|\eta B_1\|_{L^2(I)}
 \le C\bigl[(a/b)^{2+\alpha}K+\rho_i\bigr],
 \quad \rho_i\longrightarrow0.
\end{equation}
The leading constant is independent of large $N$; $\rho_i$ may depend
on fixed $N$.
\end{lemma}

\begin{proof}
Use the wider true bands in every supremum in the residual.  In collar
units the uncut regularity estimate
Equation~\eqref{eq:ricci-improvement} bounds Ricci and its scaled
derivatives on the shell of radius $r$ by $C\epsilon_i b/r$.
The dyadic sum of $r^{-1}$ is bounded, the zero-order terms carry the
summable weights $\theta_r$, and
$\theta_R/R^2\to0$.  Thus $m=o_i(1)$ uniformly.  This argument uses
the uncut estimate and is valid even where $\eta$ is small.

For the integrated assertion multiply the residual by $\eta$ and
convert a collar-unit Ricci magnitude to interval units by $b^{-2}$.
Writing $\rho=br$ for a shell's interval radius, its positive-order
terms are bounded in time $L^2$ by $V(\rho)$ from
Proposition~\ref{prop:exterior-ricci}; its zero-order term is bounded
by $\theta_r A(\rho)$.  The initial collar term is $A(b)$ up to
fixed-width changes.  Consequently it suffices to sum
\[
 A(b)+\sum_{r}\bigl[V(br)+\theta_r A(br)\bigr]
             +C\frac{\theta_R}{(Rb)^2}.
\]
The leading energy terms sum to $C(a/b)^{2+\alpha}K$.
The errors in $V$ are summable in both directions:
\[
 \sum_{br\in[cb,Cr_*]}
       \bigl[(a/(br))^\alpha+(br)^\alpha\bigr]
       \le C\bigl[(a/b)^\alpha+r_*^\alpha\bigr].
\]
The errors in $A$ are summed only with $\theta_r$, whose sum is
uniformly bounded.  In particular no constant error is multiplied
by the number of shells.  Finally
$\theta_R/(Rb)^2\to0$, because $Rb\asymp r_*>0$ and
Equation~\eqref{final:collar-weights} makes $\theta_R\to0$.
The same bounds for the fixed ample band give the assertion for
$B_1$.  Minkowski's inequality proves
Equation~\eqref{eq:collar-time-bound}.
\end{proof}

\begin{proposition}[Positive variation supplied by the true core]
\label{final:variation}
For the preceding extensions and partition,
\begin{equation}\label{eq:variation-lower}
 \sum_j\int_I\eta^2\chi_j E_j'\,dt
 \ge 2a^4K^2-a^4\bigl[CN^{-2\alpha}K^2
                              +o_i(1)(K+1)^2\bigr].
\end{equation}
\end{proposition}

\begin{proof}
The analytic collar map and its bounded derivative make $E_j$
differentiable on its patch.  The available fixed Sobolev orders
give all spatial derivatives required by the first-variation formula
in Equation~\eqref{ell:first-variation}.  On the retained true part,
\[
 \left\langle \tfrac12 Rg-\Ric,\partial_tg\right\rangle
           =2|\Ric|^2-R^2.
\]
That part contains $B_t(x_t,N_1a)$ by
Lemma~\ref{final:collar-position}, and its volume is at most $Cb^4$.
The scalar-square loss integrated in time is therefore at most
$C\epsilon_i^4b^4=o_i(a^4)$ for fixed $N$.

On the interpolated collar and attached end, use collar radius
$r\ge1$.  Equation~\eqref{eq:collar-extension}, Sobolev embedding,
and the derivative bound for the extension map imply
\[
 |\Ric(\hat g_j)|\le Cb^{-2}m(t)r^{-2-q},\qquad
 |\partial_t\hat g_j|_{\hat g_j}\le CB_1(t)r^{-q},\qquad
 d\mu_{\hat g_j}\le Cb^4r^3\,dr\,d\omega.
\]
There is no extra power of $b$ in the second estimate: normalizing
the collar metric multiplies its covariant tensor by $b^{-2}$,
whereas pulling back the original covariant Ricci tensor contributes
$b^2$ relative to its interval-unit magnitude.  The two factors
cancel.  Since $|R|\le2|\Ric|$ in dimension four, the absolute value
of the tail contribution to $E_j'$ is bounded by
\begin{equation}\label{final:tail-variation}
 Cb^2m(t)B_1(t)\int_1^\infty r^{1-2q}\,dr
       \le Cb^2m(t)B_1(t).
\end{equation}
This is the specific use of $q>1$ in the time-variation estimate.
Summing against $\chi_j$ costs no multiplicity, since
$\sum_j\chi_j=1$.  By Equation~\eqref{eq:collar-time-bound} and
Cauchy--Schwarz the integrated tail loss is at most
\[
 Cb^2\|\eta m\|_2\|\eta B_1\|_2
 \le Cb^4\bigl[(a/b)^{2+\alpha}K+\rho_i\bigr]^2.
\]
Its leading factor is exactly
\begin{equation}\label{final:tail-scaling}
 b^4(a/b)^{4+2\alpha}=a^4N^{-2\alpha}.
\end{equation}
For fixed $N$ the remaining terms are
$o_i(a^4)(K+1)^2$.  The positive core term sums to at least
$2a^4K^2$ by Equation~\eqref{final:energy-recall}.  This proves the
proposition.
\end{proof}

\subsection{Applying the static inequality to every time patch}

The lower variation bound is expressed in terms of the true root
energy. We now obtain the opposing upper bound. The only required
uniformity is over the ratios of functional value to weighted Ricci
error: a new fixed tree may be extracted from each proposed failure
of that uniformity.
\begin{proposition}[Uniform functional smallness]
\label{final:functional-smallness}
For fixed $N$ there are numbers $\zeta_i\to0$ such that, at every
time in every active patch,
\begin{equation}\label{eq:functional-littleoh}
 |\eta(t)E_j(t)|\le
 \zeta_i a^4\bigl[H_*(t)+1+a^{-2}\eta(t)m(t)\bigr].
\end{equation}
Here $H_*$ uses a fixed enlargement depending on $N$.
\end{proposition}

\begin{proof}
We first verify exactly which static hypotheses hold after an
arbitrary choice of indices, times $t_i$ with $\eta(t_i)>0$, and
active patches $j_i$.  Rescale their extensions by the length $a_i$.
Theorem~\ref{thm:degeneration} and Proposition~\ref{prop:root-scale}
extract one finite rooted tree, with the outermost nonflat entry
represented at unit scale.  Every smaller concentration is contained
inside the collar: if a concentration escaped to distance much
greater than $a$, the distance-scale limit with the root would have
two genuine tips and hence a larger nonflat entry.  For large fixed
$N$, the replacement collar and all its overlaps are uniformly
regular and leave a detecting nonflat root patch in the retained
body.

For clarity, the hypotheses of Theorem~\ref{thm:tree-inequality}
are supplied as follows.  The original compact pieces and joining
regions have the sequential identifications from
Theorem~\ref{thm:degeneration}.  On an edge with parent scale $A$
and child scale $Ae^{-2L}$, put
\[
 \tau_p=\log(A/\rho),\qquad
 \tau_c=\log(\rho/(Ae^{-2L})),\qquad \tau_p+\tau_c=2L.
\]
The physical radius is $\rho=Ae^{-\tau_p}$ on the parent half.
On the child half, $\tau_c\le L$, so
$\rho=Ae^{-2L+\tau_c}\le Ae^{-\tau_c}$.
Every parent scale is bounded by a fixed multiple of $a$.
Consequently $\rho\le C_Nae^{-D}$ on both halves, where
$D=\min(\tau_p,\tau_c)$ up to fixed end-coordinate offsets.
Equation~\eqref{eq:neck-decay} bounds the metric error by a constant
times $e^{-\mu\tau_p}+e^{-\mu\tau_c}$.  On the parent half
$\tau_p\le L$ this is at most $2e^{-\mu\tau_p}$, and on the child
half it is at most $2e^{-\mu\tau_c}$.  Thus these are exactly the
one-sided half-neck weights of the static theorem, with a possibly
smaller fixed positive exponent.  The joined groups are nontrivial
and are identified by those same annular charts.

On the new root exterior, Equation~\eqref{eq:collar-extension}
gives uniform $O(r^{-q})$ control with $1<q<2$ and sufficiently many
scaled derivatives.  Expressing it in root units changes its
constant by a factor depending on $N$, which is fixed.  Local Sobolev
compactness extracts convergence through the finite differentiability
order needed by the static theorem.  The limit is Ricci-flat on
the attached portion because $m(t_i)\to0$.  Local harmonic
coordinates and elliptic regularity make this Ricci-flat limit smooth
on its regular part; the weighted finite-order bounds are retained
in the original end coordinates.  The collar embeddings themselves
extract on smaller regular overlaps: metric and
connection bounds control their transitions with the interior
convergence charts.  Adjusting the latter by these limiting
transitions identifies the two limits on the overlaps.  Hence the
substitution changes only a regular portion and the unjoined end
of the root vertex.  It does not alter a joining end or remove the
retained nonflat patch.  This gives one fixed limiting Ricci-flat
root with the required decay $q>1$, while every other vertex is the
one furnished by the true slice convergence.  A fixed tree is
obtained only on this extracted sequence; no tree through time is
being chosen.

Let $M_i$ be a bound for the weighted, scale-neutralized Ricci error
in Equation~\eqref{eq:tree-ricci-error} on this sequence.  Choose its
internal weight exponent smaller than both the half-neck metric
exponent and the exponent in Equation~\eqref{eq:interior-weighted}.
We can take $M_i\to0$ without using $\eta$.  Indeed on a true piece
of physical radius $\rho$ in interval units the uncut bound is
\[
 \rho^2|\Ric|_{j,\rho}\le C\epsilon_i\rho.
\]
On compact pieces $\rho\le C_Na$.  At logarithmic depth $D$ in a
joining neck one has $\rho\le C_Nae^{-D}$, so multiplication by a
sufficiently small positive exponential weight still leaves a
quantity tending to zero.  The substituted region contributes at
most $C_Nm(t_i)$ with the required root-end decay.  These statements
also hold for all derivatives required in the definition of $M_i$.

Independently, Equation~\eqref{eq:interior-weighted} supplies the
stronger bound relative to the energy:
\begin{equation}\label{final:weighted-error}
 \eta(t_i)M_i
 \le C_{\mathrm{ext}}a_i^2\bigl[H_*(t_i)+\epsilon_i\bigr]
                         +C_{\mathrm{ext}}\eta(t_i)m(t_i).
\end{equation}
Here the constant may depend on this extraction and on fixed $N$.
To check the power of $a$, multiply that equation's
$(\rho/a)^2|\eta\Ric|_{j,\rho}$ by $a^2$; the result is the
scale-neutralized tensor $\rho^2|\eta\Ric|_{j,\rho}$ required in
Equation~\eqref{eq:tree-ricci-error}.  Compact retained regions are
handled by shorter regular patches, and the energy ball for $H_*$
is enlarged once to contain their short-time comparisons.  On the
collar and new end the error is $C_{\mathrm{ext}}\eta m$ by
Equation~\eqref{eq:collar-extension}.  This proves
Equation~\eqref{final:weighted-error} for the same $M_i$ obtained as
the maximum of these weighted errors.  Its uncut smallness and its
cutoff energy bound are separate assertions; neither was obtained
by division by $\eta$.

Write $E_i^{\mathrm{root}}$ for the functional in root units.
Theorem~\ref{thm:tree-inequality} now gives
$|E_i^{\mathrm{root}}|=o(M_i)$ on this extraction, and the scaling
law of Equation~\eqref{eq:energy-functional} gives
$E_{j_i}(t_i)=a_i^2E_i^{\mathrm{root}}$.  If $M_i=0$, the
zero-error scaling argument in that theorem gives
$E_i^{\mathrm{root}}=0$ directly.  Otherwise, using
Equation~\eqref{final:weighted-error} and $\epsilon_i\le1$,
\begin{equation}\label{final:ratio-test}
 \frac{|\eta(t_i)E_{j_i}(t_i)|}
 {a_i^4[H_*(t_i)+1+a_i^{-2}\eta(t_i)m(t_i)]}
 \le C_{\mathrm{ext}}\frac{|E_i^{\mathrm{root}}|}{M_i}
 \longrightarrow0.
\end{equation}

If Equation~\eqref{eq:functional-littleoh} failed uniformly, its
nonnegative ratio would be bounded below by some fixed number
along a sequence of active patches and times.  Times with $\eta=0$
have zero numerator, so that sequence has $\eta>0$.  The extraction
just performed yields Equation~\eqref{final:ratio-test} on a
subsequence, a contradiction.  The extraction-dependent constant
is fixed on that subsequence and multiplies a quantity tending to
zero.  This proves the uniform assertion without a uniform tree,
a uniform analytic family, or a uniform power-law exponent.
\end{proof}

\begin{corollary}\label{final:energy-vanishes}
Along the buffered intervals, $K_i\to0$.
\end{corollary}

\begin{proof}
Each product $\eta^2\chi_j$ has compact support in the patch where
$E_j$ is defined.  Integration by parts there gives
\[
 \sum_j\int_I\eta^2\chi_j E_j'\,dt
 =-\sum_j\int_I
       (2\eta\eta'\chi_j+\eta^2\chi_j')E_j\,dt.
\]
Absolute continuity of the partition is sufficient for this identity.
No values of different $E_j$ are equated at a switch.  By
Equations~\eqref{final:partition} and \eqref{eq:functional-littleoh},
the absolute value is at most
\begin{equation}\label{final:ibp-bound}
 C\zeta_i a^4\int_I
 [H_*+1+a^{-2}\eta m]\,[|\eta'|+\eta(1+B_1)]\,dt.
\end{equation}
For fixed $N$, Equations~\eqref{eq:maximal-ricci} and
\eqref{eq:collar-time-bound} show
\[
 \|H_*+1+a^{-2}\eta m\|_2\le C_N(K+1),\qquad
 \||\eta'|+\eta(1+B_1)\|_2\le C_N(K+1).
\]
For the first inequality use $a^{-2}=N^2b^{-2}$; fixed powers of
$N$ are allowed here.  Cauchy--Schwarz in
Equation~\eqref{final:ibp-bound} therefore bounds it by
$o_i(a^4)(K+1)^2$.

Choose and now fix $N$ so large that the leading constant in
Equation~\eqref{eq:variation-lower} satisfies $CN^{-2\alpha}<1$.
Combining the upper and lower bounds and then taking late indices
gives
\[
 K_i^2\le\gamma_i(K_i+1)^2,\qquad\gamma_i\longrightarrow0.
\]
This also proves boundedness if it was not known previously:
$K_i/(K_i+1)\le\sqrt{\gamma_i}$, and, for
$\sqrt{\gamma_i}<1$, one has
$K_i\le\sqrt{\gamma_i}/(1-\sqrt{\gamma_i})\to0$.
\end{proof}

\subsection{Small energy fixes the pointed root geometry}

The vanishing energy is an integral statement. To contradict the
prescribed change of scale, it remains to translate it into distance
preservation on a single dense set of material points and then into
preservation of the curvature-radius maximum.
\begin{lemma}[A common time-dependent Ricci bound]
\label{final:ricci-envelope}
On a slightly smaller material neighborhood of the isolated cluster,
for $0\le t\le1$,
\begin{equation}\label{eq:final-ricci-bound}
 |\Ric|(y,t)\le Q_*(t)\bigl[1+(a/l(y,t))^2\bigr],
 \qquad\|Q_*\|_{L^2([0,1])}\longrightarrow0.
\end{equation}
\end{lemma}

\begin{proof}
Inside a fixed large multiple of $a$, use the ordinary local
parabolic estimate underlying Proposition~\ref{prop:interior-ricci}
on a patch of scale a small fixed multiple of $\min(a,l(y,t))$.
The enlarged energy balls contain these patches and their prior
windows by the short-motion comparison.  Since $\eta$ is bounded
below on $[0,1]$, the estimate and
Equation~\eqref{eq:maximal-ricci} give an envelope bounded by
$C(H_*+\epsilon_i)$ times $1+(a/l)^2$.  Its time $L^2$ norm tends
to zero by Corollary~\ref{final:energy-vanishes}.

On the exterior annular range, an envelope for the magnitude on all
shells is obtained from a fixed outer shell and the sum of the
positive first derivative suprema on the intervening dyadic shells.
Indeed integrate $\nabla\Ric$ along the radial coordinate curves,
using parallel transport for tensor norms.  The controlled neck
charts give length $O(r)$ on each shell, so each shell costs its
scaled first derivative supremum.  Norms on a fixed shell vary by
the same bound along its angular paths.  Dividing the estimates for
$G=\eta\Ric$ by the fixed positive lower bound of $\eta$ is
legitimate on this interval.  Minkowski's inequality and
Proposition~\ref{prop:exterior-ricci} bound the time $L^2$ norm of
that envelope by
\[
 C\left[K+\widetilde\delta_i+
 \widetilde\delta_i\sum_{r\in[N_0a,r_0]}
                 \bigl((a/r)^\alpha+r^\alpha\bigr)\right]
 \longrightarrow0.
\]
The outer-shell value tends to zero by the uncut Ricci improvement
and its regular geometry.  Only the positive-derivative errors are
summed over all shells; the constant error in their zero-order
estimate is not summed.

Outside this range but within the smaller isolated neighborhood,
$l$ is bounded below in interval units.  Otherwise choose a sequence
with $l\to0$ there.  If its images remain in a compact regular part
of the original limiting space, smooth convergence is a
contradiction.  If they approach the isolated point, their
distance from the cluster is bounded below in interval units and
is much larger than $a$.  At that distance scale the vanishing
curvature radius and the root give two genuine tips, forcing a
nonflat entry larger than $a$.  The same argument applies when
that distance tends to infinity in interval units but tends to
zero originally, by rescaling with the distance itself.
Proposition~\ref{prop:root-scale} excludes this alternative.
Equation~\eqref{eq:ricci-improvement} thus bounds Ricci in the
remaining region by $C\epsilon_i$.  Adding these three envelopes
proves the lemma.
\end{proof}

\begin{lemma}[Good material points]
\label{final:good-points}
There is a single set of material points $\mathcal G_i$, independent
of time, whose complement has volume $o(a_i^4)$, which is $o(a_i)$-dense
at every $t\in[0,1]$ in a smaller neighborhood of the cluster.
For points $x,y\in\mathcal G_i$ in that neighborhood whose distance
at either endpoint is at most $Ca_i$,
\begin{equation}\label{final:good-distance}
 |d_1(x,y)-d_0(x,y)|=o(a_i),
\end{equation}
uniformly for each fixed $C$.
\end{lemma}

\begin{proof}
Take nested material neighborhoods compactly contained in the region
of Equation~\eqref{eq:final-ricci-bound}.  Uniform convergence of the
original distances to $d_T$ gives a fixed positive original distance
between the smaller neighborhood and the complement of the larger
one.  Thus any minimizing segment of interval length $O(a)$ with
endpoints in the smaller neighborhood remains in the region where
that equation applies.

For such a segment $\gamma:[0,L]\to M$ parametrized by arclength,
the geodesic curvature-radius estimate of
Lemma~\ref{geom:geodesic-radius} and spatial $1$-Lipschitzness of $l$
imply, on its first half,
\[
 l(\gamma(s))\ge \max\{l(x)-s,cs\}
                  \ge c'\bigl(l(x)+s\bigr).
\]
The fixed cap in that geodesic estimate is irrelevant at these
vanishing original lengths.  The second half has the analogous
bound from $y$.  Integration gives
\begin{equation}\label{final:geodesic-integral}
 \int_\gamma l^{-2}\,ds
          \le C\bigl(l(x,t)^{-1}+l(y,t)^{-1}\bigr).
\end{equation}
The length variation formula under Ricci flow, applied to minimizing
segments in both time directions, therefore yields, almost everywhere
while $d_t(x,y)\le Ca$,
\begin{equation}\label{final:distance-derivative}
 a^{-1}|\partial_t d_t(x,y)|
 \le C_C Q_*(t)
         \bigl[1+a/l(x,t)+a/l(y,t)\bigr].
\end{equation}
One may equivalently first use upper and lower one-sided length
variations; for each fixed smooth flow the distance is locally
Lipschitz in time, so these imply the almost-everywhere inequality.

The scale-sublevel covering in Theorem~\ref{thm:degeneration} gives
$\Vol_t\{l<s\}\le Cs^4$ at these scales.  Decomposing $\{l<a\}$
into the sets $2^{-k-1}a\le l<2^{-k}a$ proves
\begin{equation}\label{final:endpoint-integral}
 \int_{\{l<a\}}\frac a l\,d\mu_{g(t)}
 \le\sum_{k=0}^\infty 2^{k+1}C(2^{-k}a)^4
 \le Ca^4.
\end{equation}
Fix the material measure $d\mu_0=d\mu_{g(0)}$.  Since
$\partial_t d\mu_g=-R\,d\mu_g$ and $|R|\le C\epsilon_i^2$,
all these time-slice measures are uniformly comparable with $d\mu_0$.
Put $q_i=\|Q_*\|_{L^1([0,1])}\to0$ and define
\[
 F_i(x)=\int_0^1 Q_*(t)\frac a{l(x,t)}
                          \mathbf1_{\{l(x,t)<a\}}\,dt.
\]
Fubini and Equation~\eqref{final:endpoint-integral} give
$\int F_i\,d\mu_0\le Ca^4q_i$.  For $q_i>0$, discard the material
set $\mathcal B_i=\{F_i>\sqrt{q_i}\}$.  Then
\[
 \mu_0(\mathcal B_i)\le Ca^4\sqrt{q_i},\qquad
 \int_0^1 Q_*(t)(1+a/l(x,t))\,dt
                    \le2q_i+\sqrt{q_i}\quad(x\notin\mathcal B_i).
\]
If $q_i=0$, take the good set to be the whole manifold, up to a
null set, and the same conclusions hold.  The lower ball-volume
bound $\Vol_t B_t(x,r)\ge cr^4$ shows that every ball of radius
$C'aq_i^{1/8}$ in the smaller neighborhood meets the good set when
$C'$ is large enough.  Indeed such a ball has greater volume than
the entire bad set, uniformly in time.  Taking the complement of
$\mathcal B_i$ for $\mathcal G_i$ proves simultaneous $o(a)$ density.

For good endpoints, the time integral of the right side of
Equation~\eqref{final:distance-derivative} is $o(1)$ with a constant
depending only on a fixed distance range.  If the initial distance
is at most $Ca$, apply that inequality up to the first time it
reaches $(C+1)a$.  Its integrated variation is smaller than $a/2$
for late indices, so such a first crossing cannot occur.  The same
argument backwards in time treats a pair initially specified at
$t=1$.  This proves Equation~\eqref{final:good-distance}, and also
prevents points outside a larger bounded scaled ball from entering
a smaller one during the comparison.
\end{proof}

\begin{proposition}[The root maximum cannot change]
\label{final:root-fixed}
Let $\mathfrak a_i(t)$ denote the root scale of
Proposition~\ref{prop:root-scale}, computed in interval units.  Then
\[
 a_i^{-1}\bigl|\mathfrak a_i(1)-\mathfrak a_i(0)\bigr|
                      \longrightarrow0.
\]
\end{proposition}

\begin{proof}
Choose a good material anchor within $o(a)$ of a maximizing point
at time $0$.  Extract pointed limits at the two endpoints with
length unit $a$.  Lemma~\ref{final:good-points} identifies these
limits isometrically.  Here is the metric argument in detail.
On a fixed bounded ball choose a finite increasingly dense net of
good material labels at time $0$.  Their distance matrices at time
$1$ differ by $o(1)$ in scaled units.  The same is true of nets
chosen at time $1$, and the first-crossing argument keeps every
such label in a bounded ball at the other time.  Apply compactness
to their union and then diagonalize in the radius and the net
mesh.  The limiting correspondence preserves all pairwise
distances and is dense and onto; completing it gives a pointed
isometry.  All bounded balls used here lie near the same isolated
point of $Z$ in original units.

Every exceptional point of either limit is a genuine nontrivial
orbifold point by Theorem~\ref{thm:degeneration}, so the regular
part is determined by the metric space itself.  The isometry is
smooth there and preserves its Riemannian metric: in a sufficiently
small regular convex neighborhood, squared distances to suitable
nearby points give smooth coordinates with independent differentials;
the distance-preserving map has the corresponding smooth coordinate
expressions, as does its inverse.  Smooth convergence on these
regular regions is available at both endpoints.  In particular the
initial nonflat root patch persists in the other endpoint limit.

We next check that the curvature radius itself, rather than only
the distance metric, is preserved.  For a regular point $z$ of the
common nonflat limit $X$ define its stopped radius by
\begin{equation}\label{final:stopped-radius}
\begin{split}
 \lambda_X(z)=\sup\{r>0:\;&B_X(z,r)\text{ contains no exceptional point}\\
                       &\text{and }|\Rm_X|\le r^{-2}\text{ on }B_X(z,r)\}.
\end{split}
\end{equation}
It is finite: a regular point of nonzero curvature exists at finite
distance, so sufficiently large balls fail the curvature test.
For any regular convergent sequence $z_i\to z$ at either endpoint,
\begin{equation}\label{final:radius-convergence}
                         l(z_i,t)/a_i\longrightarrow\lambda_X(z).
\end{equation}
For the lower bound, take $r<s<\lambda_X(z)$ with $s$ admissible.
The closed $r$-ball has compact regular geometry, and its curvature
test has the strict margin $r^2|\Rm_X|\le(r/s)^2<1$.
Smooth convergence and convergence of distances imply admissibility
of every slightly smaller $r$-ball in the approximations.  Let
$r\uparrow\lambda_X(z)$.  For the upper bound, take
$r>\lambda_X(z)$ and then an intermediate radius between them.
If its ball meets an exceptional point, genuine concentration
approaches that point in the approximations and violates the
$r$-curvature test.  If it does not, its inadmissibility gives a
regular point whose curvature strictly exceeds $r^{-2}$; this
violation persists under smooth convergence.  Thus the $r$-ball
is eventually inadmissible.  The original fixed cap for $l$, divided
by the original length corresponding to $a_i$, tends to infinity,
so it imposes no further restriction here.  This proves
Equation~\eqref{final:radius-convergence}, including possible
equality cases at the limiting threshold.

The defining maximum is
\[
 \frac{\mathfrak a_i(t)}{a_i}
 =\max_{y\in\mathcal U}
   \frac{l(y,t)}{a_i}
       \psi\bigl(l(y,t)^2|\Rm|(y,t)\bigr).
\]
At a maximizer, $l/a_i$ is bounded below by a positive constant,
because $\mathfrak a_i(t)/a_i$ is bounded below and $\psi$ is bounded.
It is also bounded above by the detection and outermost-scale
property of Proposition~\ref{prop:root-scale}.  Maximizers at time
$1$ stay at bounded scaled distance from the common anchor.  Indeed
the persistent nonflat root patch supplies a nonflat entry of scale
comparable to $a_i$ at that time; if a maximizing nonflat entry
escaped from it, their distance-scale limit would have two genuine
tips and produce an entry larger than the outermost scale.  The
same bounded-range statement at time $0$ follows from the choice
of the anchor.

Consequently maximizers at either endpoint have subsequences
converging to regular points of $X$.  At any such point,
Equation~\eqref{final:radius-convergence}, smooth curvature
convergence, and continuity of $\psi$ show that their values tend
to
\[
 \lambda_X(z)\,
       \psi\bigl(\lambda_X(z)^2|\Rm_X|(z)\bigr).
\]
That same regular point can be approximated in the other endpoint
sequence, giving the same limit value there.  A subsequence
realizing any proposed strict difference between the two maxima
is therefore impossible: its maximizing point at the larger end
provides a competing point at the smaller end.  Applying this in
both directions proves the proposition.
\end{proof}

\subsection{Smooth continuation on the original manifold}

Proposition~\ref{final:root-fixed} contradicts the definite factor
change of the comparable positive root scales at the endpoints
of the intervals in Proposition~\ref{prop:buffered-intervals}.
Hence a scalar-bounded flow cannot have a concentration point.
Theorem~\ref{thm:degeneration} then gives a uniform full-curvature
bound up to the original finite time $T$.

For completeness, this yields the precise extension required in
Theorem~\ref{thm:extension}.  Curvature differentiation on a closed
manifold gives bounds for every covariant derivative of curvature
on any interval $[t_0,T)$ with $t_0>0$.  Bounded Ricci curvature and
$\partial_tg=-2\Ric$ first make $g(t)$ uniformly equivalent to
$g(t_0)$.  Integrating the connection variation and its successive
derivatives, or working in the resulting fixed controlled charts,
then gives uniform bounds and convergence for all coordinate
derivatives of $g(t)$.  Thus $g(t)$ converges smoothly to a positive
definite smooth metric $g_T$ on the original manifold $M$.
Short-time existence from $g_T$ follows from the Ricci--DeTurck
equation on this same closed smooth manifold, followed by its
diffeomorphism pullback.  The flow equation determines all time
derivatives from the spatial derivatives, so the old and new
solutions join smoothly at $T$.  Retaining the given flow for
$t<T$ and this solution for $t\ge T$ gives a smooth Ricci flow
$\widetilde g$ on $M\times[0,T+\varepsilon)$, for some
$\varepsilon>0$, with $\widetilde g(t)=g(t)$ for every $t<T$.
The argument uses arbitrary smooth initial data in real dimension
four; no change of manifold or orbifold continuation is involved.
